\subsection{Template Generation}
\label{sec:template}

$\mathrm{LoopInvTemGen}$ uses the loop body, precondition, and unchanged/non-inductive variable sets to build templates partly fixed by symbolic execution; $\mathrm{OuterLoopLLMCall}$ fills the remaining parameters (Algorithm~\supref{alg:loopinv-gen}). Prompt examples and ACSL guidance are described in \rev{\appref{sec:prompt}}.
%This template-based approach transforms LLM-based invariant synthesis into a constrained completion task,  improving the syntactic correctness of generated invariants while reducing the LLM's workload.
The six templates below cover four $\textit{SP}$-derived variable forms, one form enabled by a supplied verification goal, and a free-form \emph{sink}. The sink accommodates quantified array relations and inductive predicates that symbolic execution cannot anchor.

\begin{itemize}
    \item Unchanged variables:
For any variable $v \in uV$ and any execution path $\pi_j$ to the loop entry $p$, $v$ retains its symbolic value $\sigma_{p, \pi_j}(v)$ at the loop entry:
\[ PC_{\pi_j} \Rightarrow v = \sigma_{p, \pi_j}(v) \]


\item Loop skip conditional property: 
For a path $\pi_j$ to $p$ and a loop condition $B$, if $B$ for path $\pi_j$ is false, the loop is skipped without altering variables:
\[\rev{\neg (B \land (PC_{\pi_j} \land \rho(\sigma_{p, \pi_j}))) \Rightarrow v = \sigma_{p, \pi_j}(v)}\]
%where $\mathrm{StateForm}$ returns the logical formula specifying the values of corresponding symbolic state. % This ensures $v$ retains its initial value. 
%Both the above two types of invariants satisfy the base and preservation conditions.



\item Non-inductive variables: 
For a non-inductive variable $v$ (whose value is updated independently of its value from the previous iteration) and any path $\pi_j$ to $p$:
\[\rev{(B \land (PC_{\pi_j} \land \rho(\sigma_{p, \pi_j}))) \Rightarrow (v = \sigma_{p, \pi_j}(v)) \lor \Phi(v)}\]
where $\Phi(v)$ is the inductive invariant to be synthesized.


\item Inductive variables: 
For an inductive variable $v$ (which is neither unchanged nor non-inductive) and any path $\pi_j$ to $p$:
\[ \rev{(B \land (PC_{\pi_j} \land \rho(\sigma_{p, \pi_j}))) \Rightarrow \Phi(v)} \]

\item Verification goals:
In addition to variable-specific templates, we also introduce templates corresponding to each verification goal.
%The rationale is that, under deductive verification, many verification goals are themselves loop invariants in disguise.
For a verification goal $g$ and any path $\pi_j$ to $p$:
\[ \rev{(B \land (PC_{\pi_j} \land \rho(\sigma_{p, \pi_j}))) \Rightarrow \Phi(g)}
\]

\item Free-form sink: an unanchored slot $\Phi$ that lets the LLM propose invariants involving quantifiers and inductive predicates that the deterministic forms above cannot anchor.

\end{itemize}









Symbolic templates turn invariant synthesis into constrained completion: they provide ACSL structure and path-sensitive facts while leaving relations for the LLM. Verifier-guided refinement then retains clauses accepted by Frama-C/WP.
